% Values: PRECISE, FAITHFUL, ECONOMICAL
% Evidence planned:
% - PRECISE: a production for every bare module-level construct of the book
%   plan (module, header, imports, exposing lists and items, top-level
%   declarations, custom-type, constructor and alias bindings, type-variable
%   sequences); every restriction a decidable condition on the module alone,
%   with the recursion restriction defined by direct and delayed occurrences
% - FAITHFUL: rulings A10, A16, A17, A25, A26, A28, A30, A38, A84, A89
%   reflected
%   and cited; the default imports exactly as the survey lists them; every
%   module-level restriction of the survey ("Syntactic Restrictions",
%   "Custom Types and Aliases", "Modules") stated here or assigned by name to
%   Chapter 5 or Chapter 9
% - ECONOMICAL: the grouping of declarations left to the derived forms of
%   Appendix A; one class TopDec with abstract `;' and
%   `and'; no class of single imports; no new notation
\chapter{Syntax of Modules}\label{ch:syntax-modules}

For Modules there are further reserved words and one further symbol, but no
further special constants or identifier classes; comments and lexical analysis
are as for the Core. Every module begins with a header, which is followed by a
sequence of imports and a sequence of top-level declarations; a module makes its top-level names available
to other modules by name only. % D6
The derived forms for modules appear in
Appendix~A (Chapter~\ref{ch:appendix-derived-forms}). The \kw{infix}
declarations, \kw{effect module} headers and imports of kernel modules that
Elm admits only in its own core packages are not part of Standard Elm.
% A16, A30, A37

\section{Reserved Words}\label{sec:syntax-modules-reserved}

The following reserved words are used in Modules.
\begin{center}
\kw{as}\qquad\kw{exposing}\qquad\kw{import}\qquad\kw{module}\qquad\kw{port}\qquad\kw{type}\qquad\kw{..}
\end{center}
The words are among the reserved words of the Core
(Chapter~\ref{ch:syntax-core}), so none of them may be used as an identifier
anywhere in a program; \kw{as} is also used in patterns. The symbol \kw{..}
occurs only in exposing lists, where white space may separate it from the
parentheses that enclose it; since it is reserved, it is not an operator.
% DECIDED: `..' is a reserved symbol of the module level. Elm reads `(..)' as
% `(', `..', `)' with optional white space; no operator of the initial basis
% is `..' and a module cannot declare an operator (A16), so reserving it
% excludes no program.

The word \kw{alias} is not reserved. It has its special meaning only directly
after the reserved word \kw{type}, where a lower-case word cannot otherwise
stand, since every declared \tycon{} begins with an upper-case letter;
elsewhere \kw{alias} is a \var.

\section{Identifiers}\label{sec:syntax-modules-identifiers}

There are no further identifier classes. The class \cls{ModId} of module
identifiers is introduced with the Core, because a module identifier is the
qualifier of a long identifier \longvar, \longcon{} or \longtycon. A \modid{}
is a sequence of one or more upper-case identifiers joined by \kw{.} with no
white space between them. At this level a \modid{} has two further uses: it
names a module, in a header and in an import, and it names an alias for an
imported module, after \kw{as} in an import.

The qualifiers that may be used in a module, and the module each denotes, are
determined by the module's imports: an import \kw{import} $\modid$ \kw{as}
$\modid'$ provides the qualifier $\modid'$ and not $\modid$, and an import
without \kw{as} provides its own \modid{} (Chapter~\ref{ch:static-modules}).
A qualifier never denotes the module in which it occurs.
% A38
The identifiers declared at this level (the \var{} of a value or a port, the
\tycon{} of a custom type or an alias, the \con{} of a constructor) are never
qualified where they are declared. As in the Core, the class of each
identifier occurrence is determined by the grammatical rules which follow.

\section{Infixed Operators}\label{sec:syntax-modules-infix}

Standard Elm has no fixity declarations.
% A16
The operators that may be used in infix position in a module are those of
the initial basis (Section~\ref{sec:syntax-core-infix}). The default imports
(Section~\ref{sec:syntax-modules-default-imports}) expose every one of them
to every module, and the meaning of an operator is fixed by the initial basis,
whatever the module imports. Each has the precedence and associativity given
for it in Section~\ref{sec:syntax-core-infix}, and infixed expressions are
resolved as described there.
% Design decision 1 (audit X1): operators are fixed by the operator table of
% the initial basis; an import item (op) binds nothing.

\section{Layout}\label{sec:syntax-modules-layout}

The layout rule of Section~\ref{sec:syntax-core-layout} governs the phrases
within a top-level declaration. At the top level of a module it takes the
following form.
\begin{itemize}
\item The header, each import and each top-level declaration begins in column
  1.
\item Every other token of a header, an import or a top-level declaration lies
  in a column greater than 1. Hence a token in column 1 ends the header,
  import or declaration before it.
\item The annotation $\var\ \kw{:}\ \ty$ of a top-level value begins in column
  1, and the definition it annotates begins in column 1 on a later line, with
  only white space and comments between them; the definition repeats \var.
  % A17: the definition starts in the column of its annotation
\item A doc comment may stand directly after the header and directly before a
  top-level declaration, and nowhere else at the top level; the token that
  follows a doc comment begins in column 1. Doc comments have no meaning and
  are omitted from the grammar.
  % A10
\end{itemize}

\section{Grammar for Modules}\label{sec:syntax-modules-grammar}

The phrase classes for Modules are shown in Figure~\ref{fig:phrase-classes-modules}.
We use the variable \modm{} to range over \cls{Module}, \expitem{} to range
over \cls{ExpItem}, and the lower-case name of each other class to range over
that class. The conventions adopted in presenting the grammatical rules for
Modules are the same as for the Core. The grammatical rules are shown in
Figures~\ref{fig:grammar-modules} and~\ref{fig:grammar-top-level-declarations}.

\begin{figure}[htbp]
\centering
\begin{tabular}{@{}ll@{}}
\cls{Module}   & modules\\
\cls{Header}   & module headers\\
\cls{Imports}  & import lists\\
\cls{ExpList}  & exposing lists\\
\cls{ExpItem}  & exposing items\\[1ex]
\cls{TopDec}   & top-level declarations\\
\cls{DatBind}  & custom-type bindings\\
\cls{ConBind}  & constructor bindings\\
\cls{TypBind}  & type-alias bindings\\
\cls{TyVarSeq} & type-variable sequences\\
\end{tabular}
\caption{Modules Phrase Classes}\label{fig:phrase-classes-modules}
\end{figure}

The grammar is that of the bare language. A module as written is brought to
this form by the derived forms of Appendix~A, which gather the custom types of
a module into a single custom-type binding, and arrange its aliases and value
definitions in dependency order, the values in groups. The header is written
in every module and no derived form supplies it.
% A16, A30; book plan "Derived Forms", TopDec rows
Two symbols of the bare grammar are never written in a program. The symbol
\kw{and} joins the members of one group: the custom types of a \datbind{},
as it joins the members of a Core \valbind. The symbol \kw{;} separates two
top-level declarations, and stands for the boundary between declarations that
the layout of Section~\ref{sec:syntax-modules-layout} marks by column 1.
Neither symbol is reserved; in particular \kw{and} is an ordinary \var.

A top-level value declaration is a \valbind{} of the Core, whose bindings
carry their annotations. A port declaration consists of a \var{} and its type,
and has no definition; which types a port may have is decided by elaboration
(Section~\ref{sec:static-modules-ports}).

\begin{figure}[htbp]
\centering
\begin{grammar}
\gprod{\modm}{\header\ \imports\ \topdec}{module}
\gprod{\header}{\kw{module}\ \modid\ \kw{exposing}\ \explist}{plain}
\galt{\kw{port}\ \kw{module}\ \modid\ \kw{exposing}\ \explist}{port module}
\gprod{\imports}{\kw{import}\ \modid\ \opt{\kw{as}\ \modid'}}{import}
\galt{\quad\opt{\kw{exposing}\ \explist}\ \imports}{}
\galt{}{empty}
\gprod{\explist}{\kw{(..)}}{everything}
\galt{\kw{(}\,\expitem_1\ \kw{,}\ \cdots\ \kw{,}\ \expitem_n\,\kw{)}}{explicit}
\gnote{$(n \geq 1)$}
\gprod{\expitem}{\var}{value}
\galt{\kw{(}\opv\kw{)}}{operator}
\galt{\tycon}{type}
\galt{\tycon\,\kw{(..)}}{open type}
\end{grammar}
\caption{Grammar: Modules, Headers, Imports and Exposing Lists}\label{fig:grammar-modules}
\end{figure}

\begin{figure}[htbp]
\centering
\begin{grammar}
\gprod{\topdec}{\valbind}{value}
\galt{\kw{type}\ \datbind}{custom type}
\galt{\kw{type}\ \kw{alias}\ \typbind}{type alias}
\galt{\kw{port}\ \var\ \kw{:}\ \ty}{port}
\galt{\topdec_1\ \kw{;}\ \topdec_2}{sequential}
\gprod{\datbind}{\tycon\ \tyvarseq\ \kw{=}\ \conbind}{custom-type binding}
\galt{\qquad\opt{\kw{and}\ \datbind}}{}
\gprod{\conbind}{\con\ \ty_1\ \cdots\ \ty_n\ \opt{\kw{|}\ \conbind}}{constructor binding}
\gnote{$(n \geq 0)$}
\gprod{\typbind}{\tycon\ \tyvarseq\ \kw{=}\ \ty}{type-alias binding}
\gprod{\tyvarseq}{\tyvar_1\ \cdots\ \tyvar_k}{type parameters}
\gnote{$(k \geq 0)$}
\end{grammar}
\caption{Grammar: Top-level Declarations and Bindings}\label{fig:grammar-top-level-declarations}
\end{figure}
% DECIDED: the bare grammar keeps Elm's order `tycon tyvarseq' in datbind and
% typbind (Elm writes `type Maybe a = ...'), not the order `tyvarseq tycon'
% of the Definition of Standard ML.

In a constructor binding $\con\ \ty_1\ \cdots\ \ty_n$ each $\ty_i$ is a type
variable, a \longtycon{} without arguments, or a type enclosed in parentheses
or braces, as for the arguments of a type constructor in the Core
(Chapter~\ref{ch:syntax-core}); so a function type, or a type constructor
applied to arguments, is parenthesised when it is the argument of a
constructor. The constructor bindings of a \conbind{} are separated by
\kw{|}, which binds less tightly than any type. In a type-alias binding and a
port declaration the type extends to the end of the declaration.

\section{Default Imports}\label{sec:syntax-modules-default-imports}

The imports of every module, as elaborated (Chapter~\ref{ch:static-modules})
and evaluated (Chapter~\ref{ch:dynamic-modules}), are the default imports of
Figure~\ref{fig:default-imports} followed by the imports written in the
module. Since the default imports come first, an import written in the module
may replace a type that a default import exposes.
% A28; the operator case of A28 does not arise, since every operator is that
% of the initial basis (design decision 1).

\begin{figure}[htbp]
\centering
\begin{tabular}{@{}l@{}}
\kw{import Basics exposing (..)}\\
\kw{import Debug}\\
\kw{import List exposing ((::))}\\
\kw{import Maybe exposing (Maybe(..))}\\
\kw{import Result exposing (Result(..))}\\
\kw{import String exposing (String)}\\
\kw{import Char exposing (Char)}\\
\kw{import Tuple}\\
\kw{import Platform exposing (Program)}\\
\kw{import Platform.Cmd as Cmd exposing (Cmd)}\\
\kw{import Platform.Sub as Sub exposing (Sub)}\\
\end{tabular}
\caption{Default Imports}\label{fig:default-imports}
\end{figure}

\section{Syntactic Restrictions}\label{sec:syntax-modules-restrictions}

The following restrictions are conditions on a module as a phrase, except
those of the last item, which are conditions on a program. Each is decidable without type
information. In them, the \emph{declarations} of a \topdec{} are the value
declarations, custom type declarations, type alias declarations and port
declarations of which it is the sequence.

\begin{itemize}
\item\citelabel{res:modules-ports} In a module, the \topdec{} contains a port declaration only if the
  header is a port module header, and the \topdec{} of a module with a port
  module header contains at least one port declaration.

\item\citelabel{res:modules-header} Every module begins with a header, \kw{module} \modid{}
  \kw{exposing} \explist{} or \kw{port module} \modid{} \kw{exposing}
  \explist{}. A text that does not begin with a header is not a module of
  Standard Elm.
  % A16, A30

\item\citelabel{res:modules-import-alias} In an import, the \modid{} after \kw{as} consists of a single
  upper-case identifier.

\item\citelabel{res:modules-basics-qualifier} The qualifier \kw{Basics} is provided only by an import of the module
  \kw{Basics}: an import $\kw{import}\ \modid\ \kw{as}\ \kw{Basics}$
  requires \modid{} to be \kw{Basics}. Since no module of a program has the
  identifier of a module of the initial basis (the last item), the module
  \kw{Basics} is always that of the initial basis, and the long identifiers
  \kw{Basics.True}, \kw{Basics.False} and \kw{Basics.negate} that derived
  forms introduce (Appendix~A) always denote its bindings.
  % A89, A91; audit X18, V3

\item\citelabel{res:modules-no-operator-exposing} An item \kw{(}\opv\kw{)} occurs only in the exposing list of an import:
  a module declares no operator, so it has none to expose.
  % A16

\item\citelabel{res:modules-header-exposing} In the exposing list of a header, each item \var{} is bound by a value
  declaration or a port declaration of the module's \topdec; each item
  \tycon{} is the \tycon{} of a \datbind{} or a \typbind{} of the \topdec; and
  each item $\tycon\,\kw{(..)}$ is the \tycon{} of a \datbind, never of a
  \typbind. In particular a constructor cannot be exposed by name: an item
  \tycon{} names a type, and the constructors of a custom type are exposed
  only by $\tycon\,\kw{(..)}$ or by \kw{(..)}.

\item\citelabel{res:modules-exposing-distinct} No two items of the exposing list of a header may have the same
  identifier; thus \tycon{} and $\tycon\,\kw{(..)}$ do not both occur.

\item\citelabel{res:modules-no-top-pattern} No value declaration of a \topdec{} may contain a binding other than
  one of the form $\var\ \opt{\kw{:}\ \ty}\ \kw{=}\ \expm$; there is no
  pattern binding at the top level.

\item\citelabel{res:modules-values-distinct} No \topdec{} may bind the same \var{} twice by its value declarations
  and port declarations together.

\item\citelabel{res:modules-types-distinct} No \topdec{} may declare the same \tycon{} twice: custom types and
  aliases share one name space.

\item\citelabel{res:modules-constructors-distinct} No \topdec{} may declare the same \con{} twice. This counts the \con{}
  of each constructor binding of every \datbind, and the \tycon{} of each
  \emph{record alias}, that is, each \typbind{} $\tycon\ \tyvarseq\ \kw{=}\
  \ty$ whose \ty{} is, apart from enclosing parentheses, a record type
  $\kw{\kwlb}\,\opt{\tyrow}\,\kw{\kwrb}$ without an extension variable; a
  record alias also declares a constructor of that name
  (Chapter~\ref{ch:static-modules}).

\item\citelabel{res:modules-no-shadowing} No variable bound within an expression of a \topdec{} (a parameter of
  a lambda, a variable of a pattern, or a variable bound in a \kw{let}) may be
  a \var{} bound by a value declaration or a port declaration of that
  \topdec, wherever the declaration stands. This extends to top-level names
  the restriction on shadowing of
  Section~\ref{sec:syntax-core-restrictions}. Variables declared by imported
  modules may be bound again.
  % Audit X12: the top-level part of the shadowing restriction lives here.

\item\citelabel{res:modules-tyvarseq-distinct} No \tyvarseq{} may contain the same \tyvar{} twice.

\item\citelabel{res:modules-custom-type-tyvars} In a custom-type binding $\tycon\ \tyvarseq\ \kw{=}\ \conbind$, every
  \tyvar{} occurring in \conbind{}, including the extension variable of a
  record type, occurs in \tyvarseq. A \tyvar{} of \tyvarseq{} need not occur
  in \conbind.

\item\citelabel{res:modules-alias-tyvars} In a type-alias binding $\tycon\ \tyvarseq\ \kw{=}\ \ty$, the
  \tyvarseq{} contains exactly the type variables occurring in \ty,
  including extension variables of record types: an alias has no unused
  parameter.

\item\citelabel{res:modules-alias-not-recursive} No alias may be recursive. For the type-alias bindings of a \topdec,
  say that $\tycon$ \emph{refers to} $\tycon'$ if $\tycon'$ is declared by a
  type-alias binding of the \topdec{} and occurs, unqualified, in the type of
  the binding of $\tycon$; this relation has no cycle. The \tycon{} of a
  custom type is not related by it, so a cycle through a custom type is not
  excluded.

\item\citelabel{res:modules-value-recursion} An occurrence of a variable in an expression \expm{} is
  \emph{delayed} if it lies within a lambda expression in \expm, and
  \emph{direct} otherwise. For a value declaration \valbind{} of a \topdec,
  say that $\var \to \var'$ if \valbind{} contains a binding $\var\
  \opt{\kw{:}\ \ty}\ \kw{=}\ \expm$ and $\var'$ is bound by \valbind{} and
  has a direct occurrence in \expm. No value declaration of a \topdec{} may
  have a cycle $\var_1 \to \cdots \to \var_n \to \var_1$ $(n \geq 1)$. This
  takes the place, for the value declarations of a \topdec, of the
  restriction on recursive groups of
  Section~\ref{sec:syntax-core-restrictions}. Since a binding with
  parameters is a derived form whose equivalent form is a lambda expression
  (Appendix~A), every occurrence in its body is delayed, and a rejected cycle
  passes only through values without parameters. A delayed occurrence used
  during the evaluation of the very value it refers to passes this
  restriction; no evaluation rule applies to it
  (Chapter~\ref{ch:dynamic-modules}).
  % A26: only cycles through direct uses of parameterless values are
  % rejected at top level; A25: delayed self-reference forced at evaluation

\item\citelabel{res:modules-program} The following restrictions concern a program, and
  Chapter~\ref{ch:programs} imposes them. The \var{} of every port
  declaration of every module of a program is distinct from the \var{} of
  every other port declaration of the program and from the name of every
  manager of the platform (Chapter~\ref{ch:effects}). No module of a program
  has the identifier of a module of the initial basis. In the main module,
  \kw{main} does not occur free in the expression of its own declaration, so
  \kw{main} is not recursive, even through a delayed occurrence.
  % A84 (port names); A89 (basis module names reserved, main not recursive)
\end{itemize}

For an import, the conditions corresponding to those on the exposing list of
a header depend on what the imported module exposes: an item
$\tycon\,\kw{(..)}$ that names an alias, an item \tycon{} that names a
constructor rather than a type, and an item that the imported module does not
expose are rejected by the elaboration of imports
(Chapter~\ref{ch:static-modules}).
